\documentclass[a4paper,11pt,DIV=25]{scrartcl}

\usepackage[utf8x]{inputenc}
\usepackage[ngerman]{babel}
\usepackage{longtable}
%\usepackage{fancyhdr}
\usepackage{graphicx}
\usepackage{tabularx}
\usepackage{amsmath}
\usepackage{hyperref}
\author{Tim Reddehase}
\date{2014-09-13}

\newenvironment{oitemize}{%
  \renewcommand{\labelitemii}{$\bullet$}%
  \vspace{-15pt}
  \begin{itemize}}
  {\end{itemize}}

\begin{document}

\Huge
\begin{center}
  Ontohub loc/id\\
  \bigskip
  \includegraphics[width=0.3\textwidth]{ontohub-loc_id-logo.pdf}
\end{center}
\normalsize

\bigskip

\section{Idea}

Generally an Ontohub loc/id (locator/identifier) is just an IRI of an ontology
or one of its members (symbols, sentences, mappings, child-ontologies). However
Ontohub loc/ids are generated by the Ontohub application and assigned to an
ontology.  We try to infer them from the path of the repository, the path of
the ontology and the specific name. Additionally we ensure that this specific
IRI is actually a locator and not \textit{just} an identifier.

\section{Ontohub-Style}

Identifying Ontologies and their members in Ontohub is a hierarchical task. An
ontology belongs to a repository and a child-ontology belongs to a parent
ontology. Mappings, symbols and sentences in turn belong to a ontology. So we
should/could perfectly use the hierarchical portion of an IRI instead of the
query string. This would mean using a forward slash (\textit{/}) as separator.

Ontohub loc/ids are specific to an instance of the Ontohub-Application. However
such an instance might be reachable via multiple multiple FQDNs (fully
qualified domain name) and ports. So instead we should expect a loc/id be a
tuple consisting of the specific application instance, represented by the set
of their schema-fqdn-port tuples, and the actual identifying portion
beginning with the hierarchical forward slash (\textit{/}).

\subsection{Examples}

\begin{tabularx}{\textwidth}{p{.1\textwidth}p{.2\textwidth}p{.7\textwidth}}
  \textit{Distributed Ontology} & & \\
  \hline
  & Distributed Ontology & \url{http://ontohub.org/dol-testing/double_mapped_blendoid}\\
  & Child Ontology & \url{http://ontohub.org/dol-testing/double_mapped_blendoid/DMB-CommonSource}\\
  & Mapping & \url{http://ontohub.org/dol-testing/double_mapped_blendoid/SomeMapping}\\
  & Symbol & \url{http://ontohub.org/dol-testing/double_mapped_blendoid/DMB-CommonSource/KitchenTable}\\
  & Sentence & \url{http://ontohub.org/dol-testing/double_mapped_blendoid/DMB-CommonSource/Ax02}\\
  & & \\
  \textit{Ontology} & & \\
  \hline
  & Distributed Ontology & \url{http://ontohub.org/dol-testing/double_mapped_blendoid}\\
  & Ontology & \url{http://ontohub.org/default/pizza}\\
  & Mapping & \url{http://ontohub.org/default/pizza/SomeMapping}\\
  & Symbol & \url{http://ontohub.org/default/pizza/Veneziana}\\
  & Sentence & \url{http://ontohub.org/default/pizza/Ax02}\\
\end{tabularx}

\subsection{Conflicts in Names}

Again the loc/id of a symbol and sentence can conflict. Additionally the
mappings can conflict with symbols and sentences as well. It would be possible
to alleviate this by prepending a hierarchical portion to determine the type of
member. This could look like this:

\begin{itemize}
  \item Mapping: \url{http://ontohub.org/default/pizza/map/SomeMapping}
  \item Symbol: \url{http://ontohub.org/default/pizza/sym/Veneziana}
  \item Sentence: \url{http://ontohub.org/default/pizza/sen/Ax02}
\end{itemize}

\section{Other loc/id based IRIs}

\begin{itemize}
  \item Version: \url{http://ontohub.org/ref/2/default/pizza/map/SomeMapping}
  \item Commit: \url{http://ontohub.org/ref/def3ab/default/pizza/map/SomeMapping}
  \item Branch: \url{http://ontohub.org/ref/master/default/pizza/map/SomeMapping}
  \item Date: \url{http://ontohub.org/ref/2014-09-07/default/pizza/map/SomeMapping}
    \begin{itemize}
      \item would take the latest commit which applies to the Date range.
    \end{itemize}
  \item MMT: \url{http://ontohub.org/ref/mmt/default/pizza??SomeMapping}
    \begin{itemize}
      \item allows the use of fully compliant MMT-style
    \end{itemize}
  \item File: \url{http://ontohub.org/default/tree/some_directory/some_child_dir/Foo.txt}
    \begin{itemize}
      \item applies to HEAD commit of main branch (currently always \textit{master}
    \end{itemize}
  \item File at reference: \url{http://ontohub.org/treeref/\{ref\}/default/tree/some_directory/some_child_dir/Foo.txt}
    \begin{itemize}
      \item where \{ref\} is any of the above possible ref-types: Version, Commit, Branch or Date
    \end{itemize}
\end{itemize}

\end{document}
